\title{Byte Latent Transformer: Patches Scale Better Than Tokens}

\begin{abstract}
We present the Byte Latent Transformer (\textsc{BLT}), a novel large language model (LLM) framework operating at the byte level, that—unlike previous efforts—achieves parity with tokenization-based LLMs at scale while delivering notable enhancements in inference speed and resilience.
\textsc{BLT} processes bytes by aggregating them into adaptively determined patches, which act as the fundamental processing units throughout the model. These patches are delineated according to the predicted entropy of upcoming bytes, directing greater computational effort and representational capacity to segments characterized by higher informational complexity. We conduct the first flop-matched scaling analysis for byte-level models, extending up to 8B parameters and 4 trillion bytes of training data. Our findings underscore that effective scaling is achievable for models trained directly on raw byte sequences, dispensing with the necessity of predetermined token vocabularies. Efficiency in both training and inference is further bolstered by the strategy of dynamically grouping bytes into longer patches under conditions of low unpredictability, alongside observed qualitative gains in reasoning skills and enhanced handling of rare events. Taken together, under equivalent inference budget constraints, \textsc{BLT} demonstrates much improved scaling behavior compared to standard tokenization-based architectures, enabled by the concurrent enlargement of both patch length and model size.
\end{abstract}

\section{Introduction}
\label{section:intro}

We present the Byte Latent Transformer~(\textbf{BLT}), a novel model architecture that eschews tokenization entirely, directly leveraging raw byte sequences for learning. BLT is the first such approach to achieve parity, at scale, with traditional tokenization-based methods in terms of performance, while also offering marked gains in terms of computational efficiency and model robustness (\S\ref{sec:robustness}). In current large language model (llm) paradigms, virtually all model components are trained end-to-end with the notable exception of tokenization—a manual pre-processing technique that segments raw bytes into fixed vocabularies of tokens. The reliance on these token sets introduces inherent biases into string compression, manifesting in limitations such as domain and modality dependency~\citep{dagangetting}, elevated sensitivity to perturbations in input~(\S\ref{sec:robustness}), insufficient modeling of orthographic features~\citep{edman2024cute}, and inequities across languages~\citep{liang2023xlm,petrov2024language,limisiewicz2024myte}.

Tokenization has traditionally played a critical role, as training llms directly on raw bytes becomes computationally infeasible at scale given the extended sequence lengths~\citep{xue2022byt5}. Previous research has attempted to address this challenge by introducing more efficient variants of self-attention~\citep{el2020characterbert, clark2022canine} or by utilizing architectures that eliminate attention altogether~\citep{wang2024mambabyte}~(\S\ref{section:related_work}). Nonetheless, such strategies are mainly advantageous when training \textit{small models}. When scaling up, the bulk of computation in a Transformer arises from the sizable feed-forward network layers applied to each byte, rather than from the attention mechanism itself.

We introduce an adaptive, learnable strategy for forming \textit{patches} from sequences of bytes~(\S\ref{section:patching}) and present an architecture designed to integrate both byte-level and patch-level signals. In contrast to conventional tokenization, {\textsc BLT} dispenses with a predefined patch vocabulary. Instead, it utilizes lightweight trainable encoder and decoder components to transform arbitrary byte groupings into latent representations of patches. Our experiments demonstrate that this approach enables a \textit{more} effective use of computational resources compared to traditional models based on tokenization.

LLMs based on tokenization assign uniform computational resources to each token, regardless of the complexity of the prediction involved. This approach emphasizes performance but sacrifices computational efficiency, as the heuristics used to compress tokens do not always reflect the inherent difficulty of the prediction tasks. Our proposed architecture emphasizes the importance of adaptively distributing computational effort according to the demands of the input. For instance, predicting the ending letters of most words typically represents a straightforward, low-entropy task that does not require the full capacity of a large transformer, in contrast with the more complex task of generating the initial word in a new sentence. This principle informs the design of {\textsc BLT}'s model~(\S\ref{section:architecture}), which integrates three transformer components: two compact byte-level \textit{local models} and one expansive global \textit{latent transformer}~(\autoref{fig:arch_high_level}). 

{\textsc BLT} addresses the allocation of compute and grouping of bytes into patches by analyzing the entropy associated with predicting the next byte. This results in data segmentation that forms context-aware clusters of bytes, each containing similar levels of informational complexity.

This work introduces the first systematic study of byte-level models scaled by training flops, extending to configurations with up to 8B parameters and 4T bytes processed during training. Our results demonstrate that large models can be efficiently trained directly from byte sequences without relying on predetermined vocabulary tokenization. Notably, {\textsc BLT} achieves comparable performance to Llama 3 when evaluated under equivalent flop-controlled training regimes,\footnote{Model computational requirements are measured in Floating Point OPerations (flops).} but it offers up to a 50\% reduction in flop usage during inference~(\S\ref{section:scaling}).

In addition, we find that operating directly at the byte level yields notable advantages in modeling the infrequent and rare phenomena in datasets. Compared to tokenization-based systems, {\textsc BLT} models exhibit greater resilience to input noise and excel at character-level reasoning, as reflected in experiments on orthographic analysis, phonological inference, and tasks involving low-resource language translation~(\S\ref{sec:robustness}).

Furthermore, {\textsc BLT} architectures enable concurrent growth in both model parameter count and patch length while remaining within the same inference flop budget. Utilizing larger patch sizes decreases computation per inference, thereby allowing the saved compute to expand the global latent transformer—which is activated less frequently. Our scaling analyses at fixed inference cost~(\autoref{fig:fixed_inference_scaling}) reveal that byte-level models exhibit markedly improved scaling behavior relative to tokenized counterparts.

To summarize, our work offers the following key contributions: 1) We present {\textsc BLT}, a novel byte latent large language model that leverages dynamic compute allocation to enhance flop efficiency; 2) We demonstrate that {\textsc BLT} achieves parity with Llama 3 in terms of training flop control up to 8B parameters, and enables the possibility of achieving up to 50\% greater flop efficiency at a modest tradeoff in evaluation metric performance; 3) {\textsc BLT} models enable a novel approach to llm scaling by allowing for increases in model size without raising the inference budget; 4) We provide empirical evidence that {\textsc BLT} exhibits greater robustness to noisy input and heightened sensitivity to sub-word features that token-based llms often overlook. The training and inference code for {\textsc BLT} is publicly available at~\url{https://github.com/facebookresearch/blt}.

\section{Patching: From Individual Bytes to Groups of Bytes}
\label{section:patching}

Dividing bytes into discrete \textit{patches} enables {\textsc BLT} to flexibly adjust computational resources according to contextual requirements. Various strategies for segmenting bytes into patches are illustrated in Figure~\ref{fig:patching_types}. In precise terms, a patching function $f_p$ operates by segmenting a byte sequence $\pmb{x} = \{x_i, |i = 1, \ldots, n\}$, of length $n$, into $m < n$ patches $\pmb{p} = \{p_j | j = 1, \ldots, m\}$, using an assignment that maps each $x_i$ to either 0 or 1, where 1 designates the beginning of a new patch. For models using either tokens or patches, the principal determinant of computational expense is the number of inference steps performed by the core Transformer module. In the case of {\textsc BLT}, this corresponds to the quantity of patches required by the selected patching method to encode the dataset. Thus, the mean patch length, or simply \textit{patch size}, fundamentally governs data processing costs for both training and inference, as discussed in Section~\ref{section:flops}.

We proceed to define three distinct patching schemes: segmenting data into patches of uniform byte length~(\S\ref{section:static-patch}), partitioning at whitespace characters~(\S\ref{section:white-patch}), and forming patches adaptively using byte-level language model entropy~(\S\ref{section:dyn-patch}). We conclude with a discussion on incremental patching, as well as a comparison between patching and conventional tokenization~(\S\ref{section:bpe-patch}).

\subsection{Strided Patching Every K Bytes}
\label{section:static-patch}

A common approach for organizing bytes involves segmenting them into fixed-size patches of length $k$, as implemented in MegaByte~\citep{yu2023megabyte}. Utilizing a consistent stride simplifies both model training and deployment, while also offering a clear method to adjust the typical patch dimension, thus allowing straightforward control over computational requirements. Nevertheless, this deterministic patching approach presents notable limitations. Primarily, computational resources are not adaptively assigned to regions where they would be most effective: for example, a transformer's step $j$ might be underutilized when handling segments composed solely of whitespace in code, or conversely, may not be adequate for sections rich in content, like mathematical expressions. Furthermore, such fixed patching leads to non-uniform and context-agnostic partitioning of comparable byte sequences—for instance, resulting in inconsistent divisions of identical words.

\subsection{Space Patching}
\label{section:white-patch}

\citet{slagle2024spacebyte} introduces a straightforward but powerful enhancement to strided patching, whereby new patches are initiated following any space-like byte\footnote{
    Space-like bytes refer to those bytes which are neither Latin letters, numerical digits, nor utf-8 continuation bytes; furthermore, each patch must include at least one byte that is not space-like. }. These bytes commonly denote natural separations between linguistic elements in a wide range of languages. In the Space patching strategy, every word is explicitly modeled with an additional latent transformer step, effectively allocating greater computational resources (measured in flops) per word. This approach guarantees that words are segmented consistently across sequences, and ensures increased computational effort is devoted to positions immediately after spaces, which frequently correspond to higher uncertainty predictions. For instance, identifying the initial byte for the completion of the question ``Who composed the Magic Flute?\uline{\hspace{.6em}}'' is significantly more challenging than guessing subsequent bytes after the letter ``M'', since the set of likely first bytes is highly constrained, whereas finishing the name (i.e., ``Mozart'') becomes substantially easier. Nonetheless, the space patching technique faces difficulties in effectively supporting every language or domain, and it is restricted by its inability to adjust patch sizes flexibly. In the following, we propose a novel patching method, motivated by the observation that the initial bytes of words generally present greater prediction difficulty, while also inherently supporting variable patch sizing.

\subsection{Entropy Patching: Using Next-Byte Entropies from a Small Byte LM}
\label{section:dyn-patch}

Instead of depending on heuristic methods like whitespace, we adopt a data-centric strategy for detecting next-byte predictions with significant uncertainty. Our method, \textit{entropy patching}, leverages entropy measurements to determine the locations of patch boundaries.

We train a small byte-level auto-regressive language model on the training data for {\textsc BLT} and compute next byte entropies under the LM distribution $p_{e}$ over the byte vocabulary $\mathcal{V}$:

\begin{equation}
H(x_i) = \sum_{v \in \mathcal{V}} p_{e}(x_i=v|\pmb{x}_{<i}) \log p_{e}(x_i=v|\pmb{x}_{<i})
\end{equation}

We explore two approaches for detecting patch boundaries using the entropies $H(x_i)$. The first approach selects points whose entropy values exceed a predetermined global threshold, as shown in~\autoref{fig:patching}. The second approach flags points where the entropy is notably higher compared to the preceding value. This latter method can alternatively be viewed as locating points that disrupt the near-monotonic decline in entropy typically observed within a patch.

\begin{align*}
    \text{Global Constraint}\hspace{1cm}H(x_t)& > \theta_g\\
    \text{Approx. Monotonic Constraint}\hspace{1cm}H(x_t)& - H(x_{t-1}) > \theta_r
\end{align*}

Patch boundaries are determined through an efficient preprocessing procedure carried out at the dataloading stage. This approach contrasts with the method proposed by~\citet{nawrot-etal-2023-efficient}, which employs a classifier trained to detect entropy-based patch boundaries. In our experiments~(\S\ref{section:experiments}), we provide a comparative analysis of these two strategies for separating low and high entropy bytes.

\subsection{The Byte-Pair Encoding (BPE) Tokenizer and Incremental Patching}
\label{section:incremental-patching}
\label{section:bpe-patch}

A range of contemporary large language models, such as the Llama 3 baseline, employ subword tokenization strategies like bpe~\citep{gage1994new, sennrich-etal-2016-neural}. In this context, we define ``tokens'' as byte sequences selected from a \textit{finite} vocabulary established before model training, in contrast to ``patches,'' which are dynamically constructed sequences lacking a predetermined vocabulary. One fundamental distinction between patches and tokens is that with tokens, the model does not directly access the raw byte-level information.

An important advancement offered by {\textsc BLT} in comparison to tokenization-based architectures lies in its ability to reconfigure the balance between vocabulary size and computational demands. In conventional large language models, enlarging the vocabulary leads to an increase in the average token length, which reduces the required sequence steps but simultaneously expands the output dimensionality of the model’s final projection layer. Consequently, this inherent trade-off imposes constraints on tokenization-based models, limiting the extent to which token size and inference latency can be optimized independently. For instance, Llama 3 achieves a rise in average token length from 3.7 to 4.4 bytes, but this improvement is coupled with a fourfold expansion of its embedding matrix relative to Llama 2.

During generation, {\textsc BLT} must determine at each position in the byte sequence whether it has reached a patch boundary; this decision affects whether the Latent Transformer will be activated for additional computation. Importantly, this determination must be made solely based on the information available up to the current point in the sequence, without knowledge of subsequent bytes that are yet to be produced. Consequently, the patching mechanism must not rely on future context when segmenting the byte sequence. In formal terms, a patching strategy $f_p$ is said to exhibit incremental patching if the following holds:
$$f_p(\pmb{x}_{<i}) = f_p(\pmb{x})_{<i}$$
The bpe approach does not constitute an incremental patching scheme since identical prefixes may be tokenized differently depending on the remainder of the sequence, which violates the stated property above\footnote{Introducing a distinct delimiter token at patch boundaries can convert bpe into an incremental patching method; however, this results in an increase in byte-sequence length.}.

\section{{\textsc BLT} Architecture}
\label{section:architecture}
{\textsc BLT} consists of a sizeable global autoregressive language model that processes sequences of patch-level representations, paired with two lightweight local modules: one that transforms byte sequences into patches and another that reconstructs bytes from patch representations~(\autoref{fig:arch_high_level}).

\subsection{Latent Global Transformer Model}

\textit{The Latent Global Transformer} refers to an autoregressive transformer architecture $\mathcal{G}$ comprising $l_{\mathcal{G}}$ layers. It processes a sequence of latent input patch representations $p_j$ and outputs a corresponding sequence of patch representations $o_j$. Throughout this work, the subscript $j$ consistently identifies patch indices, while $i$ is reserved for byte indices. The global model employs a block-causal attention mask~\citep{dubey2024llama}, which enforces that each patch’s attention span is limited to itself and all preceding patches in the current document. As this model dominates the total floating point operations both during pre-training and inference, the decision of when to apply it enables flexible allocation of computational resources to input and output segments according to their respective levels of complexity.

\subsection{Local Encoder}
\label{section:local}

\textit{The Local Encoder Model}, represented by $\mathcal{E}$, is a compact model rooted in the transformer framework, distinguished by its relatively small number of layers, $l_{\mathcal{E}} << l_{\mathcal{G}}$. Its principal task is to rapidly encode an incoming stream of bytes $b_i$ into richly informative patch-level vectors $p_j$. A key modification from the conventional transformer design involves the inclusion of a cross-attention layer subsequent to each transformer layer, serving to aggregate the byte-level encodings into their corresponding patch-level representations~(\autoref{fig:crossattn}).

At the outset, the bytes $b_i$ are converted into vector representations $x_i$ using an embedding matrix in $\mathbb{R}^{256 \times h_{\mathcal{E}}}$. These byte embeddings may be supplemented with hash-embeddings to encode extra information if needed~(\S\ref{sec:hash-emb}). The model then applies multiple layers in sequence, alternately stacking transformer and cross-attention modules~(\S\ref{sec:enc-cross-attn}), yielding patch representations $p_i$, which subsequently serve as the input to the global transformer $\mathcal{G}$.

Transformer layers within $\mathcal{E}$ implement a \textit{local block causal} attention mask: each byte is allowed to attend only to the fixed-length window of $w_{\mathcal{E}}$ previous bytes. This attention window can span across patch boundaries but is prohibited from exceeding the confines of a document. The sections that follow provide an in-depth exposition of the embedding strategy and the mechanics of the cross-attention block.

\subsubsection{Encoder Hash n-gram Embeddings}
\label{sec:ngram-emb}
\label{sec:hash-emb}
An essential aspect of building resilient and informative representations at every position $i$ involves embedding details pertaining to the sequence of previous bytes. In {\textsc BLT}, this is implemented by capturing characteristics of the current byte $b_i$ alone \textit{as well as} its contribution within a byte n-gram. At each index $i$, we begin by forming byte-grams

\begin{align}
    g_{i,n}=\{b_{i-n + 1},\ldots, b_i\}
\end{align}

for every byte position $i$, and for values of $n$ ranging from three to eight.\footnote{
Byte-grams with size $n$ or larger are excluded when $i<n$.
}

Next, we propose hash $n$-gram embeddings, where each byte $n$-gram is assigned to an index in a fixed-size embedding table $E_{n}^{hash}$ using a hash function, with distinct tables for each $n \in\{3,4,5,6,7,8\}$~\citep{bai2010learning}. The computed embedding is then combined with the current byte’s embedding; the sum is subsequently normalized and supplied as input to the local encoder. The modified embedding is given by:

\begin{align}
e_i &= x_i + \sum_{n = 3,...,8}E_{n}^{hash}(\text{Hash}(g_{i,n})) \\
\text{where, Hash}(g_{i,n}) &= \text{RollPolyHash}(g_{i,n}) \% |E_{n}^{hash}|
\end{align}

Each $e_i$ is scaled by dividing by the total count of $n$-gram sizes plus one, and we implement RollPolyHash as specified in Appendix~\ref{appendix:rollpolyhash}. Section~\ref{section:ablations} investigates the role of $n$-gram hash embeddings, varying both the value of $n$ and the embedding table's capacity, to assess their influence on scaling law behavior under fixed flop budgets. Beyond employing hash-based $n$-gram embeddings, our experiments extended to frequency-based $n$-gram embeddings as well, with further explanation available in Appendix~\ref{appendix:ngrams}.

\subsubsection{Encoder Multi-Headed Cross-Attention}
\label{sec:enc-cross-attn}
Our approach is largely based on the input cross-attention component of the Perceiver architecture~\citep{jaegle2021perceiver}. The primary modification involves utilizing latent representations that map directly to variable patch representations, rather than employing a predetermined set of latent codes~(\autoref{fig:crossattn}); these representations are restricted to attending exclusively to the bytes within their respective patch. The procedure begins by generating a query vector for each patch $p_j$ through a pooling operation over the byte representations associated with $p_j$. This pooled representation then undergoes a linear transformation via $\mathcal{E}_{C} \in \mathbb{R}^{h_{\mathcal{E}} \times (h_{\mathcal{E}} \times U_{\mathcal{E}})}$, where $U_{\mathcal{E}}$ signifies the count of encoder cross-attention heads. More precisely, if $f_{\text{bytes}}(p_j)$ denotes the sequence of bytes that compose patch $p_j$, then the computation proceeds as follows:

\begin{align}
P_{0,j} &= \mathcal{E}_{C}(f_\text{bytes}((p_j)), f~ \text{is a pooling function} \\
P_l &= P_{l-1} + W_o\left(\text{softmax}\left(\frac{QK^T}{\sqrt{d_k}}\right)V\right) \\
\text{where } Q_j &= W_q(P_{l-1,j}), K_i = W_k(h_{l-1,i}), V_i = W_v(h_{l-1,i}) \\
h_l &= \text{Encoder-Transformer-Layer}_l(h_{l-1})
\end{align}

Here, $P \in \mathbb{R}^{n_p \times h_{\mathcal{G}}}$ denotes the representations of $n_p$ patches, which are input to the global model. This matrix is obtained by pooling the byte embeddings $e_i$ for each byte comprising patch $p_j$. The matrices $W_q$, $W_k$, $W_v$, and $W_o$ are projection matrices for the queries, keys, values, and outputs, with keys and values derived by projecting the byte-level representations $h_i$ from the preceding layer ($e_i$ for the initial layer). To limit attention to relevant local information, a patch-based masking mechanism is adopted so that a given query $Q_j$ can only attend to keys and values of the bytes in the same patch $j$. Since we implement multi-head attention across $Q$, $K$, and $V$, and since the patch representations are usually larger in dimension ($h_{\mathcal{G}}$) than $h_{\mathcal{E}}$, during cross-attention $P_l$ is split into several heads each of size $h_{\mathcal{E}}$, which are subsequently concatenated to reconstruct the full $h_{\mathcal{G}}$-dimensional representation. This cross-attention step includes a pre-LayerNorm operation applied to the queries, keys, and values, and does not incorporate any positional embeddings. A residual connection encloses the cross-attention component.

\subsection{Local Decoder} 

The local decoder $\mathcal{D}$ mirrors the architecture of the local encoder, utilizing a compact transformer-based framework characterized by $l_{\mathcal{D}} << l_{\mathcal{G}}$ layers. Its purpose is to transform a series of global patch embeddings $o_j$ into corresponding raw bytes $y_i$. Generating each raw byte conditioned on the preceding ones, the decoder relies on the hidden states obtained from the local encoder for the corresponding byte sequence. Processing proceeds via $l_{\mathcal{D}}$ stacked modules, each alternating between cross-attention and transformer layers. Notably, the cross-attention component is positioned before the transformer block within each layer, enabling the construction of byte-level representations from the provided patch embeddings, after which the transformer component refines these byte-wise features.

\subsubsection{Decoder Multi-headed Cross-Attention} 

In the decoder cross-attention, the roles of the queries and key/values are interchanged i.e. the byte-representations are now the queries, and the patch representations are now the key/values. The initial byte-representations for the cross-attention are initialized as the byte embeddings from the last encoder layer i.e. $h_{l_{\mathcal{E}}}$. The subsequent byte-representations for layer $l$, $d_{l,i}$ are computed as:

\begin{align}
D_0 &= h_{l_{\mathcal{E}}} \\
B_l &= D_{l-1} + W_o\left(\text{softmax}\left(\frac{QK^T}{\sqrt{d_k}}\right)V\right), \\
  &\text{where } Q_i = W_q(d_{l-1,i}), K_i = W_k(\mathcal{D}_C(o_j)), V_i = W_v(\mathcal{D}_C(o_j)) \\
  D_l &= \text{Decoder-Transformer-layer}_l(B_l)
\end{align}

Here, $W_k$ and $W_v$ denote the projection matrices used to generate keys and values, respectively; these matrices are applied to the final patch representations $o_j$ from the global model using a linear transformation followed by the split operation $\mathcal{D}_C$. The matrix $W_q$ projects the previous decoder transformer layer’s byte representations $d_{l-1}$ (or $h_{l_{\mathcal{E}}}$ for the first layer) into the query space. The output projection is conducted via $W_o$. As a result, $B \in \mathbb{R}^{h_{\mathcal{D}} \times n_b}$, with $n_b$ indicating the total number of output bytes. Subsequent decoder representations $D_l$ are then produced by applying a decoder transformer layer to $B$, the result from the cross-attention component. As done in the local encoder cross-attention mechanism, our setup leverages multi-head attention, employs pre LayerNorm, omits positional embeddings, and incorporates a residual connection that wraps around the cross-attention layer.

\section{Experimental Setup}
\label{section:experiments}
We set up systematically controlled experiments to evaluate {\textsc BLT} against tokenization-based models, ensuring that {\textsc BLT} does not benefit from potentially longer sequence lengths.

\subsection{Pre-training Datasets} In this study, all model variants are initialized using pre-training on two principal datasets: 1) The Llama 2 dataset~\citep{touvron2023llama}, which consists of 2 trillion tokens sourced from a wide array of publicly accessible domains, thoroughly curated and filtered to enhance data quality; and 2) {\textsc BLT}-1T: This is a novel compilation amounting to 1 trillion tokens, similarly assembled from public repositories and incorporating a portion of the pre-training dataset disclosed by Datacomp-LM~\citep{li2024datacomp}. The Llama 2 dataset is specifically leveraged for scaling law analyses, utilizing the guidelines on optimal token counts proposed by~\cite{dubey2024llama} to inform the most effective architectural configurations for {\textsc BLT}. Conversely, {\textsc BLT}-1T supports an entire pre-training cycle for direct comparison with Llama 3 in a range of downstream evaluations. It should be noted that neither dataset contains content derived from Meta platforms or services. In addition, when establishing baselines for models utilizing a tokenizer, we employed the Llama 3 tokenizer featuring a 128K token vocabulary, as this was observed to outperform the Llama 2 tokenizer in our benchmark experiments.

\subsection{Entropy Model} For our experiments, the entropy model employed is a byte-level language model that is trained on the identical data distribution as the complete {\textsc BLT} model. Unless specified otherwise, we utilize a transformer architecture containing 100M parameters, 14 layers, a hidden size of 512, and employ sliding window attention over sequences of 512 bytes. All other hyperparameters are aligned with those used in our local and global transformer setups. Variations in model capacity, context window size, and architectural choices are systematically explored as detailed in \autoref{section:ablations}. Notably, if the model's receptive field is sufficiently limited, the entropy model can be compactly represented as a lookup table for efficient encoding.

\subsection{Entropy Threshold and Equalizing Context Length} For models employing entropy-based patching, we determine a patching threshold designed to produce a target average \textit{patch size} over the pretraining data mixture. In the case of {\textsc BLT}, the choice of \textit{patch size} is flexible, unlike standard tokenization, which can substantially impact the context length accessible to the model. To ensure a fair comparison and prevent models with larger patch sizes from having an undue advantage, we standardize the expected number of bytes processed per batch. Consequently, models that utilize larger patch sizes are assigned shorter sequence lengths. Specifically, for Llama 2 data, we set the context to 8k bytes, while for the {\textsc BLT}-1T dataset, we increase the average context to 16k bytes, keeping the average batch size fixed at 16M bytes.

Although the mean batch size remains unchanged, the proportion of bytes to patches varies with dynamic patching approaches during data batching. To optimize performance, our implementation of {\textsc BLT} training constructs batches based on the number of patches, eliminating the need for padding within the computationally intensive latent transformer. This strategy guarantees a fixed patch count per batch. In the training phase, we pad or, if needed, truncate byte sequences to 12k and 24k bytes for the Llama 2 and {\textsc BLT}-1T datasets, respectively, mitigating the risk of memory overload due to exceptionally large patch sequences.

\subsection{Entropy Model Context}
\label{section:entropy-context}
Our empirical observations indicate that entropy patching leads to incrementally larger patches, particularly in tasks with structured formats such as multiple choice questions (as illustrated in the MMLU patching example in \autoref{fig:mmlu_patching}), where repetition is common. These differences are attributed to reduced entropy present in the repetitive segments within the entropy model context. Consequently, in our large-scale experiments using {\textsc BLT}-Entropy with a patch size of 4.5, we refresh the entropy context by introducing new lines and employ an approximate monotonicity constraint, which demonstrates greater robustness against "entropy drift" that results from fluctuations in context length. Notably, this modification only influences the entropy calculation process, and our method for determining the appropriate entropy threshold remains unchanged.

\subsection{FLOPs Estimation} 
\label{section:flops}

Our approach closely mirrors the methodology used in Chinchilla~\citep{hoffmann2022training} for calculating transformer flops, which involves tallying flops associated with feed-forward networks, qkvo projections within self-attention, as well as both attention computation and output projection. One key distinction in our analysis is the presumption that the input embedding layer leverages an efficient lookup rather than dense matrix multiplication, effectively reducing its flop contribution to zero. Consistent with prior literature, we assume that the backward pass requires twice as many flops as the forward pass.

To compute flops \textit{per byte} for {\textsc BLT} models, we add up the flops for the local encoder transformer, the global latent transformer, and the local decoder transformer, together with the cross attention blocks in the encoder and the decoder:

\begin{align}
    \text{FL}_{\text{{\textsc BLT}}} &= \frac{1}{n_p} \text{Transf. FL}(h_{\mathcal{G}}, l_{\mathcal{G}}, m=n_{ctx}/n_p,V=0)\\
   &+ \text{Transf. FL}(h_{\mathcal{E}}, l_{\mathcal{E}}, m=w_{\mathcal{E}}, V=0) \\
   &+ \text{Transf. FL}(h_{\mathcal{D}}, l_{\mathcal{D}}, m=w_{\mathcal{D}}, V=256) \\ 
    &+ k/n_p \cdot \text{Cross Attn. FL}(h_{\mathcal{E}}, l_{\mathcal{E}}, m=n_p, r=n_p/k) \\ 
   &+ \text{Cross Attn. FL}(h_{\mathcal{D}}, l_{\mathcal{D}}, m=k, r=k/n_p)
\end{align}

Here, $n_{ctx}$ represents the sequence length measured in bytes, while $n_p$ denotes the patch size. The variable $r$ specifies the ratio between queries and key/values, and $k$ indicates the proportion of the patch dimension to the byte dimension—essentially reflecting the number of local model partitions concatenated to produce a unified global model representation ($k=2$ in \autoref{fig:crossattn}). The parameter $V$ stands for the size of the vocabulary used in the output projection step, and is exclusively utilized within the local decoder. The context length, $m$, that attention mechanisms use, varies depending on whether the module operates on the byte or patch sequence. We adjust the calculation of attention FLOPs for each component based on these distinctions. Comprehensive formulations for Transformer-FLOPs and Cross-Attention FLOPs can be found in Appendix~\ref{section:flops-appendix}.

\subsection{Bits-Per-Byte Estimation} Perplexity only makes sense in the context of a fixed tokenizer as it is a measure of the uncertainty for each token. When comparing byte and token-level models, following previous work~\citep{xue2022byt5, yu2023megabyte, wang2024mambabyte}, we instead report Bits-Per-Byte (BPB), a tokenizer independent version of perplexity. Specifically:

\begin{align}
    \text{BPB}(x)&= \frac{\mathcal{L}_{CE}(\pmb{x})}{\ln(2)\cdot n_{\text{bytes}}}
\end{align}

Here, the aggregate cross-entropy loss, reflecting the uncertainty present in the data $\pmb{x}$, is scaled by dividing by both the total byte count in $\pmb{x}$ and a normalization constant.

\subsection{Transformer Architecture Hyperparameters} 

For every transformer block within {\textsc BLT}, encompassing both the local and global models, we align our design choices primarily with the Llama~3 architecture~\citep{dubey2024llama}. Specifically, the feed-forward sublayers implement the SwiGLU activation function~\citep{swiglu}, while rotary positional embeddings (RoPE)~\citep{RoPe} with $\theta=500000$~\citep{xiong2024effective} are applied exclusively in the self-attention components. RMSNorm~\citep{RMSNorm} is utilized for normalizing layers throughout the model. For self-attention operations that require fixed attention masks—namely, \textit{block causal} and \textit{fixed-window block causal} masks—Flash Attention~\citep{dao2022flashattention} is employed, setting the window size at 512 in cases where a fixed-width configuration is needed. As our cross-attention modules depend on dynamically generated, patch-specific masks, we leverage Flex Attention\footnote{\url{https://pytorch.org/blog/flexattention}}, which offers fused operations that greatly accelerate the training process.

\subsection{{\textsc BLT}-Specific Hyperparameters} To evaluate the performance of {\textsc BLT} models, we perform experiments from two perspectives: exploring scaling laws and assessing downstream application tasks. Models are trained and analyzed at various parameter scales: 400M, 1B, 2B, 4B, and 8B, with detailed architecture configurations provided in Appendix~Table~\ref{tab:arch}.

For the initialization of the first cross-attention layer within the local encoder, queries are established using max-pooling. In all {\textsc BLT} model variants, we implement $500,000$ hashes utilizing a single hash function, employing n-gram window sizes between 3 and 8. The learning rate is uniformly set to $4e-4$ for every model. This selection is informed by a hyperparameter sweep at 400M and 1B scales, scanning the interval $[1e-3, 1e-4]$, which indicated that this fixed learning rate is optimal and equivalent for both token-based and {\textsc BLT} configurations. In scaling experiments involving Llama-2 datasets, batch sizes match those prescribed by~\cite{dubey2024llama} or correspondingly scaled in data byte size. The AdamW optimizer~\citep{adamw} is adopted with parameters $\beta_1 = 0.9$, $\beta_2 = 0.95$, and $\epsilon=10^{-8}$. Training schedules incorporate a linear warm-up spanning 2000 steps, followed by cosine annealing of the learning rate down to zero; weight decay is fixed at 0.1, and global gradient norms are clipped at 1.0.

\section{Scaling Trends}
\label{section:scaling}

We offer a comprehensive overview of the scaling behavior exhibited by byte-level models to guide the continued expansion of {\textsc BLT} architectures. Our investigation into scaling dynamics is designed to overcome key shortcomings in earlier work concerning byte-level models in several respects: (a) We systematically analyze performance across the compute-optimal training regime, (b) We develop corresponding 8B parameter models using substantial training sets (reaching up to 1T tokens or 4T bytes) and assess them on a range of downstream benchmarks, and (c) We assess scaling laws under consistent inference cost constraints. A subsequent section will delve into the unique benefits inherent to byte-sequence modeling.

\subsection{Parameter Matched Compute Optimal Scaling Trends}
\label{section:parameter-matched-scaling}
We employ the Llama 2 dataset to train a series of \textit{compute-optimal} models—both bpe and {\textsc BLT} variants—spanning four model scales from 1B to 8B parameters. To assess scaling, we plot the total training flops required versus language modeling accuracy on a selected portion of the training data mixture. The bpe models are configured with the optimal model-to-data ratio as suggested by Llama~3~\citep{dubey2024llama}, ensuring a \textit{compute-optimal} allocation that should theoretically yield maximal performance for a predetermined compute budget~\citep{hoffmann2022training}, and thereby providing a solid reference point. For every bpe baseline, we additionally train a size-matched {\textsc BLT} model on identical training data, utilizing a Latent Transformer architecture that replicates the dimensions and structure of its corresponding bpe Transformer.

As depicted in~\autoref{fig:scaling-compute-opt} (right), {\textsc BLT} models consistently equal or surpass the performance of their bpe equivalents, a pattern observed across increasing model sizes and computational budgets. To our knowledge, this makes {\textsc BLT} the first byte-level Transformer architecture to demonstrate scaling behavior on par with BPE-based models when operating at compute-efficient settings. These results support our hypothesis that the ideal parameters-to-training-compute ratio for bpe remains applicable, or is at least closely similar, for {\textsc BLT}.

Both enhancements in architecture and the implementation of dynamic patching play essential roles in adhering to bpe scaling patterns. In ~\autoref{fig:scaling-compute-opt} (left), we show a comparison between models that employ space-patching and Llama 3. Our estimation of SpaceByte \citep{slagle2024spacebyte} utilizes {\textsc BLT} space-patching, omitting n-gram embeddings as well as cross-attention. While SpaceByte shows gains over Megabyte, it still does not approach the performance level of Llama 3. Meanwhile, \autoref{fig:scaling-compute-opt} (right) demonstrates the performance boosts attributable to both advancements in architectural design and dynamic patching techniques. At scale, {\textsc BLT} models achieve results comparable to top-performing tokenizer-based architectures such as Llama 3.

Additionally, we note that the selection of tokenizer significantly influences the performance of tokenizer-based models; specifically, employing the Llama-3 tokenizer yields better results than utilizing the Llama-2 tokenizer when both are applied to identical training datasets.

In summary, the {\textsc BLT} architecture exhibits behavior situated between that of Llama 2 and 3 when employed with much larger patch sizes. The average token lengths produced by the bpe tokenizers of Llama 2 and 3 are 3.7 and 4.4 bytes, respectively. Conversely, {\textsc BLT} demonstrates comparable scalability by using mean patch sizes of 6 and even 8 bytes. Since inference flops are inversely related to the average patch size, adopting an 8-byte patch size can result in approximately 50\% savings in inference computations. As both model capacity and dataset size are increased, models utilizing larger patch sizes tend to deliver improved performance. Notably, {\textsc BLT} initialized with 8-byte patches begins with significantly lower performance than bpe Llama 2 at the 1B scale but surpasses the bpe tokenizer at the 7B scale. These observations imply that larger patch sizes may prove advantageous as the model and compute scales increase, and suggest that pushing patch sizes even further might be possible as available resources grow.

\subsection{Beyond Compute Optimal Task Evaluations}
\label{section:evals}
To further investigate scaling dynamics, we train an 8B {\textsc BLT} model at a parameter-to-token ratio exceeding the compute optimal limit, utilizing the more extensive and higher-quality {\textsc BLT}-1T dataset. We evaluate model capabilities across a diverse range of established classification and generation tasks. Specifically, the evaluation covers assessments of common sense reasoning, factual knowledge, and code generation:

\paragraph{Classification tasks} examined are ARC-Easy (0-shot)~\citep{Eval_ARC}, ARC-Challenge (0-shot)~\citep{Eval_ARC}, HellaSwag (0-shot)~\citep{Eval_hellaswag}, PIQA (0-shot)~\citep{Eval_piqa}, and MMLU (5-shot)~\citep{hendrycks2020measuring}. For these, we implement a prompt-scoring approach, computing the probability distribution over candidate choices and calculate the mean accuracy for reporting results.

\paragraph{Coding related generation tasks:}  To assess the generation of Python code, we measure pass@1 on MBPP (3-shot)~\citep{Eval_MBPP} and HumanEval (0-shot)~\citep{Eval_HumanEval}, thereby evaluating the model’s ability to solve code synthesis benchmarks.

As shown in \autoref{tab:evals}, we evaluate three models, each trained using the {\textsc BLT}-1T dataset: a model leveraging the Llama 3 tokenizer with bpe,\footnote{The Llama 3 tokenizer was selected for its 128k vocabulary, which demonstrates superior performance compared to the 32k vocabulary of Llama 2.} and two forms of the {\textsc BLT} model. The first variant adopts a space-patching strategy ({\textsc BLT}-Space), while the second employs an entropy-driven patching approach ({\textsc BLT}-Entropy), enforcing an approximate monotonicity constraint and periodically refreshing the context in the entropy-based model via new lines, in accordance with the procedure described in~\autoref{section:entropy-context}. All three systems undergo training with a matched computational (flop) budget. Notably, for {\textsc BLT}-Entropy, we introduce an additional change at inference time: the entropy threshold is decreased from 0.6 to 0.1, which yields better task outcomes, though at the expense of increased inference steps.

The {\textsc BLT}-Entropy model achieves superior performance compared to the Llama 3 model in 4 of 7 evaluated tasks, despite both models being trained on an equal number of bytes. This performance gain can likely be attributed to two factors: (1) more efficient utilization of training compute made possible by dynamic patching, and (2) the explicit modeling of byte-level data rather than relying on token-based representations.

Conversely, {\textsc BLT}-Space lags behind the Llama 3 tokenizer in all tasks except one; however, due to its greater average patch size of 6 bytes, it yields a notable decrease in inference flops. For reference, both the bpe and entropy-patching models exhibit a comparable average patch size of roughly 4.5 bytes on the training set. Under an equivalent training budget, the model utilizing the larger patch size processes 30\% more data than the other two, which may further distance {\textsc BLT} from achieving compute optimality.

\subsection{Patches Scale Better Than Tokens} The {\textsc BLT} architecture allows for concurrent scaling of both model capacity and patch dimensions, all while preserving the same computational requirements for training and inference, and using an unchanged quantity of training data. This flexibility in expanding patch size is distinctive to patch-oriented approaches, enabling them to surpass the efficiency limitations inherent in token-based models that rely on static vocabularies, as detailed in Section~\ref{section:bpe-patch}. By adopting longer patches, the computation required per example is reduced, making it possible to allocate those computational resources towards expanding the global latent transformer, since its invocation becomes less frequent.

In this study, we perform a fixed inference scaling experiment to evaluate whether increasing model size while using fewer inference steps on larger input patches results in superior performance compared to smaller models requiring more steps. Utilizing models with parameter counts of 400 million and 3.6 billion based on the Llama 2 tokenizer, we identify flop-equivalent counterparts using the Llama 3 tokenizer as well as {\textsc BLT}-Entropy models that employ mean patch sizes of 6 and 8 bytes on the same training datamix (refer to~\autoref{tab:fixed-inf-params} for additional model specifications). Notably, for models with a patch size of 8, the encoder comprises 3 layers instead of just 1. Each variant undergoes training across multiple allotted training flop budgets.

\autoref{fig:fixed_inference_scaling} demonstrates that {\textsc BLT} models exhibit superior scaling behavior compared to tokenization-based models under both categories of inference flops. While BPE models provide higher performance when the training compute is limited, their advantage diminishes and is soon overtaken by {\textsc BLT} models, typically slightly past the compute-optimal point. This observation suggests that increasing investment in the pretraining stage—despite being a one-time cost—can result in a model that offers improved performance given a fixed inference compute constraint. This strategy is exemplified by 8B models such as Llama 3.1, which was trained on a volume of data that is two orders of magnitude larger than the compute-optimal amount for its parameter scale.

The point at which {\textsc BLT} surpasses token-based models in performance has moved closer to the compute-optimal region with the transition to models in the higher flop category, shifting from 3 times to 2.5 times the compute-optimal threshold. In a similar fashion, the model utilizing a patch size of 8 exhibits a more pronounced scaling trajectory in the larger flop regime, surpassing other configurations at an earlier stage. As elaborated in Section~\ref{section:parameter-matched-scaling}, increasing patch size results in performance that converges with that of BPE models as overall model size increases. We attribute this partially to the fact that the proportion of total flops consumed by the byte-level Encoder and Decoder diminishes at larger scales, likely due to their slower growth compared to the Latent Transformer. When scaling the total parameter count 20-fold, from 400M to 8B, the parameters specific to the local components of {\textsc BLT} increase by only about a factor of two. This distinction matters because increasing patch size impacts only the flop count of the Latent Transformer, leaving the byte-level components largely unaffected. Consequently, this explains why, for {\textsc BLT}-Entropy with patch size 8, the model size increased from 1.6x to 1.7x relative to Llama 2 as model scale expanded.

In conclusion, our investigation into patch-length scaling reveals that the {\textsc BLT} patch-based framework exhibits superior scaling behavior when both the patch size and the model capacity are jointly enlarged. Notably, these beneficial scaling characteristics not only persist but can further enhance as models continue to grow.

\section{Byte Modeling Improves Robustness} Additionally, we assess how {\textsc BLT} performs in terms of robustness in comparison to token-based approaches that do not access byte-level information directly. Furthermore, we introduce a method for converting existing pretrained token-based models into a byte-compatible form.
\label{sec:robustness}

\subsection{Character-Level Tasks}

An initial impetus for developing byte-level models stemmed from their enhanced tolerance to noise present at the byte level in inputs and their capacity to recognize the internal structure of tokens—a challenge for existing models that rely on tokenizers. To assess these characteristics, we conduct supplementary experiments using benchmarks specifically designed to test both input noise resilience and character-level sensitivity across English and multiple languages, covering aspects such as digits and phonemes. The findings from these experiments are summarized in Table~\ref{tab:char_tasks}.

\paragraph{Noisy Data} We generate altered variants of the benchmark classification datasets detailed in Section~\ref{section:evals} in order to assess and contrast the robustness of tokenizer-based models with that of {\textsc BLT}. To introduce textual variability, we apply five separate character-level noising techniques: (a) \textit{AntSpeak}: Rewrites the entire input in uppercase letters, separating each character with spaces. (b) \textit{Drop}: Deletes 10\% of the characters at random positions within the input text. (c) \textit{RandomCase}: Randomly assigns uppercase or lowercase to each character so that 50\% are uppercase and the remaining 50\% lowercase. (d) \textit{Repeat}: Selects 20\% of the text's characters and duplicates them, repeating each character up to four times. (e) \textit{UpperCase}: Changes every character in the text to uppercase. For each evaluation scenario, we apply these noise types separately to the prompt, the completion, or both, and present the mean performance across all cases. Table \ref{tab:char_tasks} summarizes the performance results on noisy HellaSwag~\citep{Eval_hellaswag}, demonstrating that {\textsc BLT} consistently achieves superior robustness compared to tokenizer-based models, with an average lead of 8 points over a model trained on identical data, and it even surpasses the Llama 3.1 model, which was trained on a substantially larger dataset.

\paragraph{Phonology - Grapheme-to-Phoneme (G2P)} We evaluate how effectively {\textsc BLT} converts sequences of graphemes—i.e., the character representations of words—into their corresponding pronunciations, represented as phonemes. Table \ref{tab:char_tasks} reports the outcomes of the G2P task under a 5-shot setting, utilizing the Phonology Bench~\citep{suvarna2024phonologybench}. The results indicate that {\textsc BLT} achieves superior performance relative to the baseline Llama 3 1T tokenizer-based model on this phonological mapping task.

\paragraph{CUTE}
To evaluate character-level comprehension, we test {\textsc BLT} using the CUTE benchmark~\citep{edman2024cute}. This benchmark encompasses a variety of tasks that are generally grouped into three main areas: comprehension of composition, recognition of orthographic similarity, and sequence manipulation skills. For many models relying on tokenizers, these tasks prove challenging; such models tend to have some awareness of token spellings, yet often fail to leverage this capability for manipulating textual data. According to Table~\ref{tab:char_tasks}, {\textsc BLT}-Entropy surpasses both BPE-based Llama 3 variants by over 25 points on this suite of tasks. Notably, our model excels at character manipulation, attaining 99.9\% accuracy in both spelling-related tasks. These substantial gains occur even though {\textsc BLT} has been trained on 16 times less data compared to Llama 3.1, underscoring that acquiring character-level knowledge is particularly difficult for BPE approaches. Figure \ref{fig:cute_examples} provides examples where Llama 3 models based on BPE encounter difficulties but the {\textsc BLT} model achieves strong performance. The only areas where BPE models hold an advantage are word deletion and insertion; although such operations may be less straightforward for byte-level architectures, the performance discrepancy is modest, suggesting that constructing word-level understanding from characters may be more feasible than the inverse. For all tasks, we adopt the original evaluation methodology and prompts provided by Huggingface. It is worth noting that BPE-based models could potentially see improvements with more extensive prompt engineering.

\paragraph{Low Resource Machine Translation} We assess the performance of {\textsc BLT} on translation tasks both to and from 21 lower resource languages spanning six prominent language families, each utilizing a diverse array of scripts, based on the FLORES-101 benchmark~\citep{goyal2022flores}. SentencePiece BLEU scores are presented in Table~\ref{tab:flores}. The findings reveal that {\textsc BLT} surpasses a model utilizing the Llama 3 tokenizer, exhibiting an overall improvement of 2 BLEU points for translations into English, and a 0.5-point gain for translations from English. While {\textsc BLT} matches or marginally exceeds Llama 3 in well-represented language pairs, it delivers more substantial improvements across many language pairs in the context of lower-resource language families. This highlights the strength of byte modeling in enabling better generalization to rare or diverse byte sequences.

\subsection{Training {\textsc BLT} with Llama 3 Initialization}
\label{sec:distillation}
We investigate an approach where {\textsc BLT} models build upon pre-existing pre-trained models that utilize tokenizers, thereby facilitating faster and more robust convergence during training. Specifically, we initialize the global transformer parameters of {\textsc BLT} with those from a pre-trained Llama 3.1 checkpoint. During subsequent training, the global transformer's weights are updated using a learning rate set to one-tenth of that used for the local encoder and decoder modules. Training proceeds for the same number of optimization steps as typically used for Llama 3, and results are compared with a standard {\textsc BLT} in Table~\ref{tab:distill}. The results clearly demonstrate that {\textsc BLT} initialized from Llama 3.1 substantially outperforms both the Llama 3 baseline and the {\textsc BLT} baseline, while matching their computational budget. Additionally, even in comparison to our {\textsc BLT}-Entropy variant (see Table~\ref{tab:evals}), which was trained on a substantially larger corpus (1T tokens), {\textsc BLT} initialized from Llama 3.1 exhibits stronger performance on the MMLU benchmark, highlighting its potential to dramatically cut down on training flops.

This approach can be interpreted as converting models that rely on tokenizers into ones that do not, thereby transforming a pre-trained LLaMA 3.1 model into a {\textsc BLT} model. For a thorough evaluation, Table \ref{tab:distill} presents results for the original LLaMA 3.1, trained on 15T tokens, and compares its performance with that of its {\textsc BLT} counterpart. While the transformed model shows only minor reductions in performance on MMLU and HumanEval benchmarks, more pronounced declines are observed across additional tasks. These outcomes indicate that additional research is required to more effectively harness the capabilities of the pre-trained model, with particular attention to refining data mixtures and tuning hyperparameters.

\section{Ablations and Discussion}
\label{section:ablations}
Here, we present ablation studies that support the architectural decisions underlying {\textsc BLT} as well as the selection of the patching method and associated hyper-parameters used for the {\textsc BLT} model with 8B parameters, trained on the {\textsc BLT}-1T dataset.

\paragraph{Entropy Model Hyper-parameters} To investigate how the size of the entropy model and the length of the context window influence scaling dynamics, we experiment with byte-level entropy transformer models, varying their sizes from 1m to 100m parameters and using context window lengths ranging from 64 up to 512. We generate scaling law plots of bits-per-byte (bpb) as a function of training flops, utilizing our $400m$ and $1b$ {\textsc BLT} models trained on the Llama-2 dataset; these plots are displayed in \autoref{fig:entropy_ablation}. Our analysis reveals a positive relationship between scaling performance and both the model size and the window length, although gains tend to level off beyond 50m parameters.

\paragraph{Types of Patching}

We conduct an ablation study of the four distinct patching approaches outlined in Section~\ref{section:patching}: 1) Strided Patching with strides set to 4 and 6, 2) Patching at whitespace boundaries, 3) Patching via BPE Tokenizer utilizing the Llama 3 tokenizer, and 4) Entropy-driven patching leveraging a compact byte-level language model.

Although dynamic patching shortens the actual sequence length, we normalize sequence lengths across different patching methods to ensure that the contextual window remains consistent. This means that each model processes, on average, an equivalent number of bytes per sequence during both training and inference, which eliminates the possibility of confounds arising from variations in context size. The findings from these controlled comparisons are illustrated in Figure~\ref{fig:scaling-compute-opt}. All dynamic patching strategies examined demonstrate better performance compared to static patching, with space patching performing nearly as well as dynamic entropy-based patching.

\autoref{tab:evals_ablation} displays comparative benchmark results for {\textsc BLT} models utilizing tokenizer-based approaches, space patching, and entropy-based patching. These models were trained on the Llama 2 dataset for an optimized number of training steps~\citep{dubey2024llama}. While space patching offers a more straightforward approach, eliminating the need to execute an entropy model dynamically during training, our results demonstrate that the improvements associated with entropy-based patching in scaling behavior~(see Section~\ref{section:scaling}) also translate to better performance on downstream benchmarks.\footnote{The outcomes for space patching are sourced from prior experiments that did not incorporate cross-attention, but analogous patterns persist even when cross-attention is included.}

\paragraph{Cross-Attention} 
\autoref{tab:crossattn_ablations_1b} presents results from experiments removing and modifying cross-attention at different stages within both the encoder and decoder components of {\textsc BLT}. Regarding encoder cross-attention, we consider several approaches to query initialization: (1) employing a shared learned embedding for all global states, (2) utilizing a hash embedding corresponding to the bytes within the patch, and (3) deriving the query from a pooled version of the encoder's hidden states over the patch bytes at the current layer.

Our results indicate that integrating cross-attention within the \textit{decoder} yields the greatest benefit. Applying cross-attention in the encoder leads to minor gains, and these are observed solely when the queries are initialized via pooling. Moreover, the advantages of cross-attention are most pronounced on the Common-Crawl dataset, with the effect further amplified for larger patch sizes.

\paragraph{n-gram Hash Embeddings} 

We conduct an ablation study using n-gram hash embedding vocabulary sizes of 0, 100K, 200K, and 400K, with corresponding outcomes shown in Table~\ref{tab:ngram_ablations_1b}. The findings suggest that incorporating hash embeddings improves model performance across all evaluated domains, with a more pronounced impact on Wikipedia and Github datasets (yielding a 0.04 bpb increase versus 0.01 bpb after 15k training steps at the 8B scale). At the same 8B model scale, reducing the number of hashes from 500K to 300K results in a performance change of approximately 0.001 bpb at 15k steps, demonstrating that hash embeddings are crucial for enabling {\textsc BLT} models to achieve performance comparable to traditional tokenizer-based approaches. Nevertheless, increasing the number of hashes beyond 300K yields minimal additional benefit. Moreover, the performance enhancements obtained through hash embeddings are largely additive when combined with cross-attention, as these methods yield improvements across distinct datasets.

\paragraph{Local Model Hyperparamaters} Table \ref{tab:local_layers} presents an ablation study examining different configurations for the depth of the local encoder and decoder. When used in conjunction with hash n-gram embeddings, {\textsc BLT} demonstrates strong performance when the encoder is kept minimal—with only a single layer—while a more substantial decoder yields better results.

\section{Related Work}
\label{section:related_work}

\textbf{Character-Level RNNs:} Since the advent of neural models~\citep{sutskever2011generating,mikolov2012subword,graves2013generating}, Character Language Modeling has gained popularity because it naturally handles out-of-vocabulary words, eliminating the need for explicit back-off strategies. \cite{kim2016character} develop a model that processes only characters on the input side, employing convolutional and highway networks whose outputs serve as inputs to LSTM-based RNNs. This approach demonstrated comparable results to state-of-the-art RNN language models of the time for English, while offering notable performance gains for languages with rich morphological structures, highlighting another critical benefit of character-level LLMs. In a related vein, \cite{kenter2018byte} apply byte-level LSTM models for machine comprehension tasks and report higher accuracy than word-level models, particularly for Turkish and Russian, both morphologically complex languages. Additionally, \cite{zhang2015character} apply convolutional models based on characters to various classification tasks, outperforming their word-level counterparts on several benchmarks. Expanding on hierarchical representations, \citet{chung2019hierarchical} propose hierarchical LSTM models equipped with boundary detectors to uncover latent hierarchical structures in text, further enhancing the performance of character-level language models. In machine translation, ByteNet~\citep{kalchbrenner2016neural} utilizes CNN layers on characters rather than relying on attention mechanisms.

\textbf{Character-Level Transformers:} The advent of attention-based transformer architectures~\citep{vaswani2017attention}, combined with the use of subword tokenization strategies~\citep{sennrich-etal-2016-neural}, brought substantial gains in the effectiveness of neural models across various language modeling and benchmark evaluations. Nevertheless, utilizing word or subword tokens imposes an implicit inductive bias regarding the abstraction layer these models handle. Seeking to integrate the achievements of transformer models with early positive outcomes in character-level language modeling, \citet{al2019character} leveraged deeply stacked transformers, employing auxiliary loss functions to train these models such that they surpassed previous LSTM-based character language models. Despite these advances, a notable disparity between character-level and word-level LLM performance remained. Similarly, GPT-2~\citep{radfordgpt2} noted that, when evaluated on extensive datasets such as the 1 Billion Word Benchmark, byte-level language models still lagged behind their word-level counterparts in terms of competitiveness.

Although \cite{Choe2019BridgingTG} show that transformer-based language models operating at the byte level can surpass subword-level LLMs with a comparable number of parameters, they note that such models require substantially greater computational resources and extended training durations. In a related vein, \cite{el2020characterbert} introduce CharFormer, a variant of BERT that constructs word representations by applying convolutional layers to character embeddings, demonstrating notable gains in the biomedical domain—but at the cost of significant compute expenditure. Building further, \cite{clark2022canine} propose CANINE, an encoder-only model containing 150M parameters that processes raw character sequences. Central to CANINE is a deep stack of transformers akin to our global model, alongside a hybrid approach employing local transformers and strided convolution to efficiently downsample the sequence of characters; this architecture achieves better performance than a token-level encoder-only baseline (mBERT) on a range of multilingual downstream tasks. The ByT5 model~\citep{xue2022byt5} explores strategies for byte-level encoder-decoder models that fully omit patching operations. Although ByT5 displayed higher resilience to noisy input and competitive performance relative to token-based models using only a quarter as much data, the absence of patching necessitated resource-intensive attention over individual bytes, leading to substantial computational overhead. Modeling language directly at the byte level rather than with subword units elongates input sequences, exacerbating the challenge of scaling such models efficiently. Leveraging the recently introduced Mamba Architecture~\citep{gu2023mamba}, which supports fixed-size memory across vast context windows, \cite{wang2024mambabyte} develop a byte-level Mamba model also free from patching operations. This model surpasses byte-level transformers in a flop-matched regime, at the 350M parameter scale, in terms of bits-per-byte performance on multiple datasets.

\textbf{Patching-based approaches:} Leveraging patching strategies can significantly reduce the computational cost, measured in flops, commonly associated with byte-level LLMs, without necessarily sacrificing model performance. Several studies have shown promising results for these methods at modest scales in terms of both model parameters and training data. For instance, \cite{nawrot-etal-2022-hierarchical} explore static patching techniques involving downsampling and upsampling operations and introduce the hourglass transformer architecture, which achieves superior results compared to other byte-level models at the 150M parameter level. Building on this, \cite{nawrot-etal-2023-efficient} present enhancements via dynamic patching strategies, such as an end-to-end trainable boundary predictor, a version of the boundary predictor guided by specific tokenizers, and an entropy-based approach similar to {\textsc BLT}. Their experiments show that such dynamic patching methods surpass the performance of conventional transformers of the time when evaluated on language modeling at the 40M parameter scale using approximately 400M training tokens. Meanwhile, \cite{lester2024training} study model training on sequences encoded with arithmetic coding, thereby attaining compression levels exceeding those possible with BPE. Utilizing an equal-info windows technique, they demonstrate improved results relative to byte-level baselines on language modeling tasks; however, their methods fall short compared to subword-based baselines.

This study is primarily motivated by MegaByte~\citep{yu2023megabyte}, a causal decoder-only LLM employing a predetermined, static approach for patching and concatenating representations, which facilitates the mapping from bytes to patches. On the decoder's side, MegaByte implements a local model that reconverts these patches into bytes. Their results indicate that, at a scale of 1B parameters and using a corpus comprising roughly 400B bytes, MegaByte is capable of performing comparably with models that utilize traditional tokenization. In our comprehensive experiments, we include ablations of MegaByte and observe that static patching methods consistently trail behind leading compute-optimized tokenizer-based models within settings where floating-point operations are controlled. We further illustrate the capacity of {\textsc BLT} to close this performance gap. Likewise, \cite{slagle2024spacebyte} identify the same limitations of MegaByte and propose adaptations such as patching based on whitespaces and other similar byte types, along with the integration of a local encoder. Their findings reveal performance gains over transformer models relying on tokenization in certain domains—specifically, Github and arXiv—at the 1B parameter level under compute constraints. We include results with their approach in our analysis and demonstrate that additional advances in model architecture are necessary to further advance byte-level models and achieve parity with top-performing token-based systems like Llama 3.

\section{Limitations and Future Work}

In this study, we adopt the training durations identified as optimal for Llama~3~\citep{dubey2024llama} when selecting model architectures. Nevertheless, the referenced scaling laws originate from analyses conducted on BPE-level transformers, and their direct application may produce less-than-ideal ratios between data and parameter counts in the context of {\textsc BLT}. Determining scaling laws specifically tailored to {\textsc BLT} is left for future investigation, which might uncover even more advantageous scaling behaviors for our approach. Furthermore, most experiments in this work utilize model scales up to 1B parameters; as we extend our exploration to larger models, such as those with 8B parameters or more, different architectural optima may emerge, possibly resulting in enhanced performance at those larger scales.

Current transformer codebases and libraries are primarily optimized for architectures that rely on tokenizers. Although our study includes theoretical comparisons matched for floating-point operations and incorporates some optimized components—like FlexAttention—for handling non-standard transformer layers, our current implementations may still lag behind tokenizer-based models regarding actual runtime performance, and additional improvements could further enhance their efficiency.

Whereas {\textsc BLT} relies on a distinct entropy model for patching that is trained independently, an alternative avenue for future research would be to integrate the learning of the patching model directly within an end-to-end training framework. In Section \ref{sec:distillation}, we detail preliminary experiments that suggest the feasibility of converting tokenizer-based architectures—such as Llama 3, which are trained on over 10T tokens—into byte-based models by initializing the global transformer with their pre-trained weights and subsequently freezing those parameters. Pursuing this line of investigation may reveal strategies that simultaneously maintain the advantages offered by byte-based modeling and potentially achieve superior performance relative to these tokenizer-based architectures, all without necessitating training from the ground up.

\section{Conclusion}
\label{section:conclusion}

In this work, we introduce the Byte Latent Transformer (\textbf{BLT}), an innovative architecture that moves beyond the standard reliance on fixed-vocabulary tokenization commonly used in large language models. {\textsc BLT} employs an adaptive, learnable strategy for assembling bytes into patches, enabling the model to allocate compute in response to the intrinsic complexity of input data. This adaptability yields notable enhancements in both model robustness and computational efficiency. Through a comprehensive scaling analysis, we show that {\textsc BLT} models achieve performance parity with tokenization-centric models such as Llama 3 at scales of up to 8B parameters and 4T bytes. Furthermore, BLT can offer up to a 50\% decrease in inference flops with minimal sacrifices in downstream evaluation performance. Additionally, {\textsc BLT} introduces the possibility to expand both the model and the patch sizes simultaneously without exceeding a prescribed inference cost, thus facilitating more flexible scaling behaviors suitable for realistic computational settings. By operating directly at the byte level, {\textsc BLT} enhances the model’s capacity to interpret rare or atypical data, strengthening its resilience to noisy inputs and enriching its representation of sub-word information. Collectively, these advancements position {\textsc BLT} as a compelling and versatile alternative to established tokenization-based methodologies, offering a robust and resource-efficient path forward for next-generation language models.

\end{document}